% Generated draft: requires mathematical and source review.
\documentclass[11pt]{article}
\usepackage[margin=25mm]{geometry}
\usepackage{fontspec}
\setmainfont{Latin Modern Roman}
\usepackage{amsmath,amssymb,mathtools,unicode-math}
\setmathfont{Latin Modern Math}
\usepackage{enumitem,needspace,xurl,hyperref}
\setlength{\parindent}{0pt}
\setlength{\parskip}{5pt}
\setlength{\emergencystretch}{3em}
\newcommand{\N}{\mathbb{N}}
\newcommand{\Z}{\mathbb{Z}}
\newcommand{\Q}{\mathbb{Q}}
\newcommand{\R}{\mathbb{R}}
\newcommand{\C}{\mathbb{C}}
\newcommand{\F}{\mathbb{F}}
\newcommand{\A}{\mathbb{A}}
\newcommand{\PP}{\mathbb{P}}
\newcommand{\E}{\mathbb{E}}
\newcommand{\eps}{\varepsilon}
\newcommand{\dd}{\,\mathrm{d}}
\begin{document}
\section*{UnsolvedMath ID: 30002280}
\textbf{Source problem code:} catalogue entry
\subsection{Definitions and mathematical statement}
A real-valued function \(f\) on an open convex set \(U\subseteq\R^m\) is called a difference-of-convex (DC) function if there exist finite convex functions \(g,h:U\to\R\) such that \(f=g-h\). This definition explains the terminology; it does not specify how the boundaries in the conjecture must be represented by such functions.

\textbf{Source conjecture, with unresolved qualifications.} Kinematic formulas hold for all sets in the intended class whose boundaries are described by differences of convex functions and for which normal cycles can be constructed.

The supplied record does not define the intended class, the ambient space or dimension, the boundary representations, or the normal cycles required. It also does not identify the kinematic identities being asserted, their participating sets, curvature quantities, transformation group, integration measure, or normalization. Accordingly, the assertion above preserves the source conjecture but is not a fully specified mathematical proposition. In particular, the DC-function definition alone supplies no missing boundary regularity, compactness, or integrability hypotheses. It is also unresolved whether existence of a normal cycle is an additional hypothesis or an asserted property of the intended boundary class.
\subsection{Short English statement}
The conjecture asks whether kinematic formulas extend to an intended class of sets with boundaries described by differences of convex functions and with constructible normal cycles. The class and the formulas are unspecified, so a precise quantified formulation requires further source clarification.
\subsection{Sources}
\begin{itemize}
\item Local UnsolvedMath registry, permanent ID \(30002280\), source problem code \texttt{catalogue entry}, title: \emph{Kinematic Formulas for DC-Boundary Sets}. The supplied record belongs to the Oberwolfach Reports Open Problems collection and gives proposed year 2013; no individual proposer is identified.
\item Supplied report citation: \emph{Integral Geometry and its Applications} (2013), Oberwolfach Reports, pp.\ 295--342, DOI 10.4171/owr/2013/05. Source URL and supplied literature-source link, labeled ``original Oberwolfach report'': \url{https://doi.org/10.4171/owr/2013/05}. The background describes Oberwolfach Reports as documenting mathematical workshops and their open-problem discussions. No specific page or contributor for this conjecture is supplied.
\item Literature assessment dated \texttt{2026-08-21}: supplied dated curation, not independently verified. The supplied status is ``open.'' The assessment says that the kinematic-formula claim is explicitly conjectural, that the intended class definition is omitted, and that no general proof was found in the supplied literature triage. The record asks that time-sensitive status and cited claims be verified independently.
\item The registry assigns statement status \texttt{corrected\_verified} and reports a manual semantic, typographic, and LaTeX review describing the cleanup as self-contained and meaning-preserving. This is supplied review metadata; it does not resolve the missing mathematical specifications or establish the conjecture.
\end{itemize}
\end{document}
